\documentclass[11pt]{article}
\usepackage{mathblatt}


\begin{document}
\pruefheftstil{Prozent \textperiodcentered{} Grundwert}{Prüfungsheft P10 \textperiodcentered{} Fokus}
\noindent{\Large\bfseries Prozent \textperiodcentered{} Grundwert}\hfill{\small\color{mbgrau}Prüfungsheft P10 \textperiodcentered{} Fokus}\par\medskip
\pfabschnitt{Prozent}
\pfstufe[sgrundwert]{Grundwert}
\kommtdran{in 2 von 13 Jahren (2023, 2025) \textperiodcentered{} zuletzt 2025 \textperiodcentered{} 1 BE}
\begin{pfaufgabe}{1.}{1 BE \textperiodcentered{} P10 2025 · 1a}
Ein T-Shirt ist um 6 € billiger geworden. Dieser Nachlass entspricht 20 \% des ursprünglichen Preises. Geben Sie den ursprünglichen Preis an.
\rechenplatz{3}
\fusshilfe{1: 30 €, Tipp: Prozentsatz}
\end{pfaufgabe}
\pfweiterekopf
\pfpaar{\pfkopfzeile{2.}{eigene Aufgabe}
$36\,\%$ einer Menge sind $81$ kg. Wie viel kg ist die ganze Menge?
\par\noindent\hfill \leerfeld[kg]
\fusshilfe{2: $225$ kg}}{\pfkopfzeile{3.}{eigene Aufgabe}
$21$ Kinder fahren mit dem Rad. Das sind $75\,\%$ der Klasse. Wie viele Kinder hat die Klasse?
\par\noindent\hfill \leerfeld
\fusshilfe{3: $28$ Kinder}}
\pfpaar{\pfkopfzeile{4.}{eigene Aufgabe}
$64\,\%$ einer Strecke sind $52$ m. Wie lang ist die ganze Strecke?
\par\noindent\hfill \leerfeld[m]
\fusshilfe{4: $81{,}25$ m}}{\pfkopfzeile{5.}{eigene Aufgabe}
Im Tank sind noch $12$ l. Das sind $20\,\%$ der Tankfüllung. Wie viele Liter passen in den Tank?
\par\noindent\hfill \leerfeld[l]
\fusshilfe{5: $60$ l}}
\pfpaar{\pfkopfzeile{6.}{1 BE \textperiodcentered{} P10 2023 · 1b}
Mia gewinnt bei einem Wettbewerb Geld. Sie gibt 25 \% des Gewinns an ihre Eltern ab; das sind 200 €. Geben Sie an, wie hoch der gesamte Gewinn war.
\fusshilfe{6: 800 €, Tipp: Prozentsatz}}{\pfkopfzeile{7.}{eigene Aufgabe nach P10 2025}
Beim Kauf einer Mütze spart Mira $4\,€$, das sind $25\,\%$ des alten Preises. Wie viel kostete die Mütze vorher?
\par\noindent\hfill \leerfeld[€]
\fusshilfe{7: $16\,€$}}
\pfpaar{\pfkopfzeile{8.}{eigene Aufgabe nach P10 2025}
Bei einer Konzertkarte beträgt der Rabatt $15\,€$, das sind $30\,\%$ des alten Preises. Wie hoch war der alte Preis?
\par\noindent\hfill \leerfeld[€]
\fusshilfe{8: $50\,€$}}{\pfkopfzeile{9.}{eigene Aufgabe nach P10 2023}
Lisa bekommt $40\,\%$ einer Erbschaft, das sind $6\,000\,€$. Wie groß ist die ganze Erbschaft?
\par\noindent\hfill \leerfeld[€]
\fusshilfe{9: $15\,000\,€$}}
\pfpaar{\pfkopfzeile{10.}{eigene Aufgabe}
Tims Handy zeigt $28\,\%$ Akku. Laut Anzeige reicht das noch für $7$ Stunden Musik. Wie lange hält ein voller Akku? Reicht ein voller Akku für eine Busfahrt von $22$ Stunden?
\fusshilfe{10: Ja}}{\pfkopfzeile{11.}{eigene Aufgabe}
Die Klassenfahrt ist zu $70\,\%$ bezahlt. Das sind $4\,200\,€$. Wie teuer ist die Klassenfahrt insgesamt? Wie viel Geld fehlt noch?
\fusshilfe{11: $1\,800\,€$}}
\pfpaar{\pfkopfzeile{12.}{eigene Aufgabe}
Nach einem trockenen Sommer sind in einem Stausee noch $18$ Mio.\,m$^3$ Wasser. Das sind $45\,\%$ der Füllmenge. Wie viel Wasser fasst der See? Wie viel fehlt noch, bis er voll ist?
\fusshilfe{12: $22$ Mio.\,m$^3$}}{}
\end{document}